\documentclass{resume}

\usepackage[left=0.4in,top=0.4in,right=0.4in,bottom=0.4in]{geometry}
\usepackage{fontawesome5}
\newcommand{\tab}[1]{\hspace{.2667\textwidth}\rlap{#1}}
\newcommand{\itab}[1]{\hspace{0em}\rlap{#1}}
\name{Kevin Lee}
\address{+1 530-376-9732 \\ Davis, CA}
\address{\faEnvelope\ \href{mailto:yhylee@ucdavis.edu}{yhylee@ucdavis.edu} \quad
\faLinkedin\ \href{https://www.linkedin.com/in/kevin-lee-b78448388/}{linkedin.com/in/kevin-lee-b78448388} \quad
\faGithub\ \href{https://github.com/section9-us}{github.com/section9-us}}

\begin{document}
\small

\begin{rSection}{PROFESSIONAL SUMMARY}
AI-focused Software Engineer with 10+ years of experience applying software engineering, security analytics, and data-driven investigation to critical infrastructure and public-sector environments. Currently pursuing an M.S. in Computer Science focused on machine learning, LLMs, generative AI, RAG, model evaluation, and MLOps. Strong at turning complex technical data, threat evidence, and domain-specific workflows into reliable AI-assisted systems and analytical solutions.
\end{rSection}

\begin{rSection}{EXPERIENCE}

\textbf{Software Engineer, Threat Management Division} \hfill Jan 2022 -- Present\\
Ministry of Science and ICT \hfill \textit{Sejong, Republic of Korea}
\begin{itemize}
    \itemsep -4pt {}
    \item Produced 500+ CTI cases per year by analyzing malware, attack evidence, and APT behaviors, creating structured technical intelligence that can support AI-assisted classification, retrieval, detection, and incident response workflows.
\end{itemize}

\textbf{Software Engineer, Security Planning Division} \hfill May 2018 -- Jan 2022\\
Ministry of Science and ICT \hfill \textit{Sejong, Republic of Korea}
\begin{itemize}
    \itemsep -4pt {}
    \item Conducted AI-assisted technical security assessments and vulnerability analysis for 50+ public institutions annually (200+ total), applying analytical methods to identify patterns, prioritize risk, and support remediation for critical infrastructure systems.
\end{itemize}

\textbf{Software Engineer, Reactor Design Development Division} \hfill Nov 2014 -- Mar 2018\\
KEPCO E\&C \hfill \textit{Gimcheon, Republic of Korea}
\begin{itemize}
    \itemsep -4pt {}
    \item Developed reactor design software and an engineering database for nuclear plant workflows over 3+ years, building practical experience with domain-specific data modeling, software automation, and secure engineering systems.
\end{itemize}

\end{rSection}

\begin{rSection}{SKILLS}
\begin{tabular}{ @{} >{\bfseries}l @{\hspace{2ex}} p{0.77\textwidth} }
Programming & Python, Java, C/C++, C\#, JavaScript, HTML/CSS, SQL, Spring MVC \\
AI/ML & Machine Learning, LLMs, Generative AI, Prompt Engineering, Fine-Tuning, RAG, Agentic Systems, Model Evaluation, and MLOps \\
Security Data & Malware Analysis, CTI, APT Analysis, Incident Response, Vulnerability Analysis, Threat Modeling, and Security Analytics \\
Tools & PyTorch, TensorFlow, Scikit-learn, Hugging Face, MySQL, PostgreSQL, Oracle, Docker, Burp Suite, Nmap, Wireshark, and IDA Pro \\
Languages & Korean (Native), English (Fluent) \\
\end{tabular}\\
\end{rSection}

\begin{rSection}{EDUCATION}

{\bf Master of Science in Computer Science}, University of California, Davis (Current GPA: 4.0/4.0) \hfill {Expected Jun 2027}\\
Korean Government Long-term Fellowship for Overseas Study, Ministry of Science and ICT

{\bf Bachelor of Science in Computer Software}, Kwangwoon University (GPA: 3.32/4.0) \hfill {Feb 2015}

\end{rSection}

\begin{rSection}{CERTIFICATIONS \& TRAINING}
\textbf{Certifications:} Engineer Information Processing, HRDK (2013) \quad SQL Developer, Korea Data Agency (2014) \\
\textbf{Training:} Locked Shields, NATO CCDCOE (2023) \quad Advanced Web Application Exploits, Black Hat Training USA (2021)
\end{rSection}

\end{document}
